\documentclass[aspectratio=169,11pt]{beamer}
% Shared CMPSC 480 Beamer preamble.
\usepackage[T1]{fontenc}
\usepackage[utf8]{inputenc}
\usepackage{lmodern}
\usepackage{amsmath,amssymb,mathtools}
\usepackage{booktabs,array,tabularx}
\usepackage{xcolor}
\usepackage{listings}
\usepackage{tikz}
\usetikzlibrary{arrows.meta,positioning,shapes.geometric}

% Ignore only negligible TeX box diagnostics that are below visual tolerance.
% Larger layout defects still appear in the build log and in rendered QA.
\vfuzz=3pt
\hbadness=2000

\definecolor{courseink}{HTML}{111111}
\definecolor{coursegray}{HTML}{EDEDED}
\definecolor{courserule}{HTML}{B8BCC4}
\definecolor{courseblue}{HTML}{3D8DFF}
\definecolor{courselightblue}{HTML}{D0EDFA}
\definecolor{coursemuted}{HTML}{5C6470}
\definecolor{coursegreen}{HTML}{0B7A53}
\definecolor{coursered}{HTML}{B3261E}

\usetheme{default}
\usefonttheme{professionalfonts}
\setbeamertemplate{navigation symbols}{}
\setbeamersize{text margin left=0.65cm,text margin right=0.65cm}

\setbeamercolor{normal text}{fg=courseink,bg=white}
\setbeamercolor{frametitle}{fg=courseink,bg=white}
\setbeamercolor{title}{fg=courseink,bg=white}
\setbeamercolor{subtitle}{fg=coursemuted,bg=white}
\setbeamercolor{structure}{fg=courseblue}
\setbeamercolor{alerted text}{fg=coursered}
\setbeamercolor{block title}{fg=courseink,bg=courselightblue}
\setbeamercolor{block body}{fg=courseink,bg=coursegray}
\setbeamercolor{example text}{fg=coursegreen}

\setbeamerfont{title}{size=\fontsize{25}{29}\selectfont,series=\bfseries}
\setbeamerfont{subtitle}{size=\fontsize{13}{16}\selectfont}
\setbeamerfont{frametitle}{size=\fontsize{19}{22}\selectfont,series=\bfseries}
\setbeamerfont{normal text}{size=\fontsize{11}{14}\selectfont}
\setbeamerfont{footline}{size=\fontsize{7.5}{9}\selectfont}

\setbeamertemplate{frametitle}{%
  \nointerlineskip
  % Use content-driven height so a long assessment title can wrap without
  % being clipped by a fixed-height title bar.
  \begin{beamercolorbox}[wd=\paperwidth,sep=0.17cm,leftskip=0.48cm,rightskip=0.48cm]{frametitle}
    \insertframetitle
  \end{beamercolorbox}
  \vspace{-0.08cm}
  {\color{courserule}\rule{\textwidth}{0.35pt}}
}

\setbeamertemplate{footline}{%
  \leavevmode
  \begin{beamercolorbox}[wd=\paperwidth,ht=2.6ex,dp=1.1ex,leftskip=.65cm,rightskip=.65cm]{author in head/foot}
    \color{coursemuted}\insertshorttitle\hfill\insertframenumber/\inserttotalframenumber
  \end{beamercolorbox}
}

\setbeamertemplate{itemize item}{\color{courseblue}\small$\blacktriangleright$}
\setbeamertemplate{itemize subitem}{\color{courseblue}\scriptsize$\bullet$}
\setbeamertemplate{enumerate item}{\color{courseblue}\insertenumlabel.}

\lstdefinestyle{coursepython}{
  language=Python,
  basicstyle=\ttfamily\fontsize{8.5}{10}\selectfont,
  keywordstyle=\color{courseblue}\bfseries,
  commentstyle=\color{coursemuted}\itshape,
  stringstyle=\color{coursegreen},
  showstringspaces=false,
  columns=fullflexible,
  keepspaces=true,
  frame=single,
  rulecolor=\color{courserule},
  backgroundcolor=\color{coursegray},
  breaklines=true,
  tabsize=4,
  xleftmargin=0.15cm,
  xrightmargin=0.15cm,
  framexleftmargin=0.15cm,
  framexrightmargin=0.15cm
}
\lstset{style=coursepython}

\newcommand{\points}[1]{\hfill{\color{courseblue}\textbf{[#1 points]}}}
\newcommand{\deliverable}[1]{\textbf{\color{courseblue}#1}}
\newcommand{\code}[1]{\texttt{#1}}
\newcommand{\keyidea}[1]{%
  \begin{beamercolorbox}[rounded=false,sep=0.55em,wd=\linewidth]{block body}
    \textbf{Key idea.} #1
  \end{beamercolorbox}
}
\newcommand{\answerlabel}{\textbf{\color{coursegreen}Answer.}\ }
\newcommand{\gradinglabel}{\textbf{\color{courseblue}Grading guidance.}\ }

\newenvironment{questionblock}[1]
  {\begin{block}{#1}}
  {\end{block}}

\newenvironment{solutionblock}[1]
  {\begin{block}{#1}}
  {\end{block}}

\AtBeginSection[]{
  \begin{frame}
    \vfill
    {\usebeamerfont{title}\color{courseink}\insertsectionhead\par}
    \vspace{0.4cm}
    {\color{courseblue}\rule{0.32\paperwidth}{3pt}}
    \vfill
  \end{frame}
}


% CMPSC480_TOTAL_POINTS: 10
% CMPSC480_ITEM_IDS: 1,2,3,4
\title[Week 2 Assignment Questions]{Week 2 Assignment: A Unified Search Engine}
\subtitle{BFS, DFS, UCS, and A-star behind one tested graph-search core}
\author{CMPSC 480 - Student Question Deck}
\date{}

\begin{document}


\begin{frame}
  \titlepage
\end{frame}

\begin{frame}{Assessment overview}
\begin{columns}[T,onlytextwidth]
\begin{column}{0.62\textwidth}
\Large Implement and compare four search strategies while preserving one domain contract.

\vspace{0.35cm}
\normalsize
Coverage: Weeks 1-2: problem formulation, graph search, costs, heuristics, and testing
\end{column}
\begin{column}{0.33\textwidth}
\begin{block}{Points}10 total\end{block}
\begin{block}{Expected time}6-8 hours\end{block}
\begin{block}{Submission}Repository plus PDF report\end{block}
\end{column}
\end{columns}
\end{frame}

\begin{frame}{Learning objectives}
\begin{itemize}
  \item Implement BFS, DFS, UCS, and A-star through a shared frontier abstraction.
  \item Use state keys and best-cost bookkeeping correctly in cyclic graphs.
  \item Justify admissibility and consistency of a grid heuristic.
  \item Compare strategies with solution-cost and search-effort measurements.
  \item Design edge-case tests that distinguish algorithmic guarantees.
\end{itemize}
\end{frame}

\begin{frame}{Requirements checklist}
\small
\begin{itemize}
  \item Use Python 3.11 and NumPy only unless the prompt explicitly permits another package.
  \item Preserve the published interfaces and make all tie-breaking deterministic.
  \item State assumptions, validate inputs, and handle empty, terminal, and unreachable cases.
  \item Include focused tests whose names explain the defect each test would expose.
  \item Report measured evidence; do not replace experimental results with unsupported claims.
\end{itemize}
\end{frame}

\begin{frame}{Task 1 - Shared search contract (2 points)}
\small
Define the problem, node, result, and frontier interfaces before implementing a strategy.

\vspace{0.2cm}
\begin{enumerate}
\item[\textbf{a.}] Specify the required fields and methods of SearchProblem, Node, and SearchResult. \hfill{\color{courseblue}\textbf{[1 points]}}
\item[\textbf{b.}] Define an immutable state-key policy and list two input validation rules. \hfill{\color{courseblue}\textbf{[1 points]}}
\end{enumerate}
\end{frame}

\begin{frame}{Task 2 - BFS and DFS policies (2 points)}
\small
Implement breadth-first and depth-first graph search by changing only frontier removal order.

\vspace{0.2cm}
\begin{enumerate}
\item[\textbf{a.}] Explain the exact frontier and explored-set behavior for BFS and DFS. \hfill{\color{courseblue}\textbf{[1 points]}}
\item[\textbf{b.}] Give one test that separates BFS from DFS and one cyclic test that proves termination. \hfill{\color{courseblue}\textbf{[1 points]}}
\end{enumerate}
\end{frame}

\begin{frame}{Task 3 - UCS and A-star - part 1 (4 points)}
\small
Add cost-sensitive search and demonstrate that cheaper rediscovery replaces a worse frontier path.

\vspace{0.2cm}
\begin{enumerate}
\item[\textbf{a.}] State the priority used by UCS and A-star and specify when the goal test occurs. \hfill{\color{courseblue}\textbf{[1 points]}}
\item[\textbf{b.}] Trace a graph where A is first reached with cost 5 and later through B with cost 2. \hfill{\color{courseblue}\textbf{[1 points]}}
\item[\textbf{c.}] Define Manhattan distance for a four-neighbor unit-cost grid and prove admissibility. \hfill{\color{courseblue}\textbf{[1 points]}}
\end{enumerate}
\end{frame}

\begin{frame}{Task 3 - UCS and A-star - part 2 (4 points)}
\small
Add cost-sensitive search and demonstrate that cheaper rediscovery replaces a worse frontier path.

\vspace{0.2cm}
\begin{enumerate}
\item[\textbf{d.}] Prove consistency of Manhattan distance for each legal edge. \hfill{\color{courseblue}\textbf{[1 points]}}
\end{enumerate}
\end{frame}

\begin{frame}{Task 4 - Experiments and defensive tests (2 points)}
\small
Evaluate all strategies on common instances and explain differences using theory.

\vspace{0.2cm}
\begin{enumerate}
\item[\textbf{a.}] Report path cost, generated nodes, expanded nodes, and maximum frontier for at least three maps. \hfill{\color{courseblue}\textbf{[1 points]}}
\item[\textbf{b.}] Include no-solution, single-cell, cheaper-rediscovery, and scale tests and state the bug each detects. \hfill{\color{courseblue}\textbf{[1 points]}}
\end{enumerate}
\end{frame}

\begin{frame}{Edge cases and defensive programming}
\small
\begin{itemize}
  \item The initial state already satisfies the goal test.
  \item The graph contains a reachable cycle but no goal.
  \item A state is rediscovered through a strictly cheaper path.
  \item Two frontier entries have equal priority and require deterministic ties.
  \item The heuristic returns zero everywhere, reducing A-star to UCS.
\end{itemize}
\end{frame}

\begin{frame}{Submission instructions}
\small
\begin{itemize}
  \item Submit source, tests, and a README containing the exact test command.
  \item Include a PDF result table and a short explanation of every surprising measurement.
  \item Tag the final commit and ensure a clean clone reproduces the tests.
\end{itemize}
\vspace{0.2cm}
\alert{Do not submit credentials, generated caches, or unapproved libraries.}
\end{frame}

\begin{frame}{Grading rubric}
\scriptsize
\begin{tabularx}{\textwidth}{>{\bfseries}p{0.28\textwidth} p{0.08\textwidth} X}
\toprule
Category & Points & Required evidence \\
\midrule
Shared search contract & 2 & Typed interfaces, validation, and a deterministic result schema. \\
BFS and DFS policies & 2 & Working FIFO and LIFO policies plus trace tests. \\
UCS and A-star & 4 & Priority-frontier implementation, heuristic, and correctness traces. \\
Experiments and defensive tests & 2 & Reproducible result table, tests, and concise analysis. \\
\bottomrule
\end{tabularx}
\end{frame}

\begin{frame}{Reflection prompt}
In 150-200 words:

Identify one implementation choice that affected a theoretical guarantee. Cite a failing or passing test and explain what you changed.
\end{frame}

\end{document}
